\section{Introduction}

A world model deployed online has to notice when the world stops behaving as it
expects. The usual signal is the model's own prediction error: when the next
observation does not match the prediction, something has changed. This is the
basis of surprise-driven monitoring in model-based control, and it has the
appealing property of being free --- the model already computes it.

But prediction error is a scalar summary of a high-dimensional internal state,
and it only moves once the mismatch has reached the output. A natural
conjecture is that the internal state moves first: that the \emph{geometry} of
the representation --- how sharply its trajectory bends, how far it has drifted
from its recent history, how its local covariance has rotated --- registers the
disturbance before the loss does. If so, a monitor reading the geometry would
buy reaction time at no extra model cost, which is exactly the currency that
matters for a controller that has to stop or re-plan.

The conjecture is plausible enough to be worth testing and vague enough to be
easy to confirm by accident. This paper does the test properly.

\paragraph{What makes this hard to get right.}
Three failure modes sit between a plausible experiment and a trustworthy one,
and each quietly favours a predetermined answer.

\begin{enumerate}
\item \textbf{The baseline can cheat.} V-JEPA 2's reference implementation runs
one encoder over the whole clip and then uses token masks to split context from
target. The context representations have therefore already attended to the very
frames being predicted. The resulting loss is partly self-fulfilling, which
makes the baseline look \emph{stronger} than it is --- biasing against the
hypothesis under test. We take context from a separate encoder pass whose final
tubelet is overwritten, and verify by assertion that changing the predicted
frames leaves the prediction bit-identical (Section~\ref{sec:method}).

\item \textbf{Thresholds can leak.} An alarm threshold fitted on the same clips
used to measure false alarms reports an optimistic operating point. We split
control clips into disjoint calibration and evaluation sets, and cross-fit the
learned baseline within the calibration set so that even its threshold is
honest.

\item \textbf{The headline metric can be degenerate.} Time-to-alarm is
right-censored: clips that never alarm have no observed value. Once fewer than
half the clips alarm inside the window, the median is pinned at the censoring
value \emph{for every monitor}, and the comparison returns ``no difference''
regardless of the truth. We registered the median as the primary metric anyway,
because it was specified in advance --- and registered a censoring-robust
secondary alongside it, before running anything, precisely so that a degenerate
primary would be visible rather than silently null.
\end{enumerate}

\paragraph{A fourth trap, specific to this question.}
Even a perfectly executed comparison can be vacuous. If the prediction-error
baseline does not detect the anomalies at all --- if its detection rate sits at
its own false-alarm rate --- then ``geometry warns earlier than prediction
error'' is true in the same way that anything is earlier than never. It says
nothing about whether geometry is an early warning signal. We therefore report a
\emph{lift} for every monitor (detection rate minus held-out false-alarm rate),
flag any baseline that fires at chance, and mark a passing verdict as
\textsc{vacuous} when the probe that passed only beat such a baseline. We know
of no prior work on representation-based anomaly detection that reports this
check, and we think it should be standard.

\paragraph{Contributions.}
\begin{itemize}
\item A pre-registered, leakage-free protocol for the question ``does internal
geometry warn before prediction error?'', with causality enforced structurally
(monitors are streaming objects that cannot reference the future) and verified
adversarially (inverting all frames after $t$ leaves every score unchanged).
\item A benchmark of four anomaly families injected at a fixed frame into a
controlled physical scene, with matched pre-injection histories so that
comparisons are paired at the clip level.
\item An evaluation that reports time-to-detection as survival analysis ---
Kaplan--Meier curves, log-rank, paired Wilcoxon, and a restricted mean
time-to-alarm with a paired cluster bootstrap --- under two threshold regimes,
with multiplicity control across all comparisons.
\item A model-agnostic adapter interface: the causal loop, probes, thresholds
and statistics are shared, so a new architecture is added in roughly twenty-five
lines and is held to the same contract by a parameterised test suite. This is
what makes a cross-architecture claim a comparison rather than a collection of
incomparable pipelines.
\item The registered result on frozen V-JEPA 2: \lgwVerdict.
\end{itemize}

\paragraph{On reporting a negative result.}
The protocol commits in advance to a decision gate, and to writing up a negative
result if no probe clears it. We regard that commitment as the main reason to
trust whatever the gate returns. The harness is designed to make a null answer
as publishable, and as informative, as a positive one.
